\section{An infinite family attaining the position bound}
\label{app:two-prime}
Let $n=2p$ with $p$ an odd prime. Put $s=n(n-3)/2$,
$a=(n-1)(n-2)/2$, $d=s+a=n^2-3n+1$, and $t=s/d<1/2$.
Choose the cyclic addition table of $\mathbb Z_n$ with probability
$w=t(n-1)/p$ and the following square otherwise. Both probabilities are positive.
In the second square, slots and symbols belong to $\mathbb F_p\times\{0,1\}$.
The $2p$ rows are the following maps, with $i\in\mathbb F_p$ and
the first coordinate calculated modulo $p$:
\[
\begin{aligned}
 A_i(x,c)&=(x+(-1)^c i,c),\\
 B_i(x,c)&=(x+(-1)^c i+c,1-c).
\end{aligned}
\]
Each row permutes the symbols. For a fixed slot, the two row families
range once over opposite fibers, proving the Latin property. Independently
permute slots and symbols uniformly after choosing the square.

\paragraph{Relative rows.}
Condition on two distinct uniformly drawn rows. Within either fiber
of row indices, the relative row translates both copies of $\mathbb F_p$
by nonzero increments, so its cycle type is $(p,p)$.
Across fibers, $A_i^{-1}B_j$ interchanges the copies and adds
$(-1)^c(i+j)+c$ to the first coordinate. Squaring adds one on both
copies, so the relative row is a single $2p$-cycle. In the cyclic
table a nonzero translation is equally likely to be any of the
$n-1$ increments. Exactly one increment yields $p$ transpositions,
and $p-1$ yield each long-cycle type. Thus in the mixture the probability
of type $2^p$ is $w/(n-1)=t/p$, the expected number of transpositions
in the cycle decomposition is $t$, and the mean sign is $-1/(n-1)$.
Indeed the probability of an even relative row in either table is
$(p-1)/(2p-1)$.

\paragraph{The low-dimensional components.}
By Eq.~\eqref{eq:character-certificate}, it suffices to bound the
off-diagonal mean character divided by degree by $1/d$ on every
nontrivial irreducible. The symmetric and skew dyad components have
dimensions $s,a$ and mean characters $t,1-t$, giving equality in
both. On a derangement with $u$ transpositions their characters are
$u,1-u$. Every row pair with $u>0$ has negative sign. The two
sign twists therefore have mean characters $-t,t-1/(n-1)$, giving
ratios at most $1/d$ since $t<1-t$. The sign character and standard
character have negative means. The sign twist of the standard character
has ratio $1/(n-1)^2<1/d$.

\paragraph{All remaining components, uniformly in the prime.}
For $n\ge256$, the low-degree classification of Rasala
\citep{rasala} gives degree $f\ge n(n-1)(n-5)/6\ge n^3/7$
after removing the trivial representation and the seven nontrivial
components treated above. On an $n$-cycle the absolute character is
at most one by the Murnaghan--Nakayama rule. On $(p,p)$ it is at
most two: in a beta-set encoding each removable $p$-hook moves one
bead to a vacant position $p$ places lower. All initially available
moves have disjoint endpoints and can be executed together. Three
moves would remove $3p$ boxes from a shape of size $2p$. After
the first move the absolute character on the remaining $p$-cycle is
at most one. On $2^p$ the Fomin--Lulov inequality \citep{fominlulov}
gives absolute character at most
\[
 \frac{2^p p!}{\sqrt{(2p)!}}\sqrt f\ \le\ (2n)^{1/4}\sqrt f .
\]
Here $\binom{2p}{p}\ge4^p/\sqrt{4p}$ follows by induction from $p=1$.
Since $t/p<1/n$, the absolute off-diagonal ratio on all these components
is at most
\[
 \frac{2}{f}+\frac{(2n)^{1/4}}{n\sqrt f}
 \le \frac{1}{n^2}\left(\frac{14}{n}+
          \frac{2^{1/4}\sqrt7}{n^{1/4}}\right)<\frac1{n^2}<\frac1d.
\]
The parenthesis is less than one already at 256 and decreases.
Adding back the diagonal gives largest centered eigenvalue
$1/n+(n-1)/(nd)=L_n$. Theorem~1 supplies the matching correctness
lower bound. This proves simultaneous optimality for every prime
$p\ge131$, without a separate search at each size.
